\section{二人零和世界：Minimax、Exploitability 与 ``稳健但不一定最赚''}

\subsection{从扑克问题切入：两种 ``最好'' 的冲突}
Brown 先抛出课堂投票问题：谁是更好的扑克玩家？
\begin{itemize}
    \item 选项 A：对任何对手长期 head-to-head 都不亏（稳健）；
    \item 选项 B：一年内总盈利最高（收益最大化）。
\end{itemize}
这两个选项在国际象棋/围棋里往往重合，但在扑克里可能不重合。A 对应 Minimax/Nash 语义，B 更接近对群体分布的 exploit 策略。

\begin{knowledgebox}{术语对齐}
\begin{itemize}
    \item \textbf{Minimax Equilibrium}：在最坏对手下保证不亏损的策略集合。
    \item \textbf{Exploitability}：相对最优反制（best response）的可被剥削程度。
    \item \textbf{Population Best Response}：针对某个对手分布最大化收益的策略。
\end{itemize}
\end{knowledgebox}

\subsection{Exploitability 的工程意义}
课程中用 Rock-Paper-Scissors 举例。若策略总是出 Rock，则对手最佳反制是总出 Paper，exploitability 极高。若等概率随机化三种动作，则 exploitability 为 0。

这个例子对应到线上 Agent 部署非常现实：当模型面向海量用户暴露接口，策略弱点会被群体发现并复用，因此 ``最坏情况稳健性'' 不能忽略。

\begin{importantbox}{可部署系统的底线}
当策略会长期暴露在开放环境中，先压低 exploitability，再追求额外收益，通常比反过来更稳妥。
\end{importantbox}

\subsection{Minimax 为什么在扑克里仍然 ``能赢''}
Brown 给出关键解释：在复杂不完美信息游戏中，Minimax 不等于 ``只能打平''。因为对手会犯负期望错误（negative EV actions），你即便只守住均衡，也会在长期统计上获利。

这与课堂原句一致：\textit{``As your opponents make mistakes, you profit.''} 该命题在德扑职业策略里几乎是共识。

\begin{warningbox}{误区：把 Minimax 误解成保守策略}
Minimax 是风险边界，不是被动保守。它提供的是 ``不被系统性击穿'' 的地板，而非 ``永不利用对手'' 的天花板。
\end{warningbox}

\subsection{Sound Self-Play 的条件}
Brown 反复强调：在二人零和前提下，sound self-play 才有收敛保证。这里 sound 不是口号，至少包含：
\begin{itemize}
    \item 足够探索（sufficient exploration）；
    \item 对策略混合分布的正确逼近；
    \item 训练过程不会系统性塌陷到可被反制的局部策略。
\end{itemize}

在理论极限（无限容量和算力）下，这一性质很强：无需人类示范数据，也能收敛到不亏策略。

\subsection{本章小结}
二人零和提供了一个少见的 ``理论与工程同向'' 场景：Minimax 定义清晰、可用 exploitability 量化、可由 sound self-play 逼近。但这组性质不应被外推到所有 Agent 任务。

